\begin{monsheet}
\addcontentsline{toc}{section}{\ExamTocProblemLabel}

\begin{daimon}{1}
次の各問いに答えよ。\hfill（配点 20点）
\begin{shomon}
  \item 不定方程式 $5x+8y=200$ を満たす正の整数の組 $(x, y)$ の個数を求めよ。また，そのような組 $(x, y)$ のうち，$x^{2}+y^{2}$ の値を最小にする組 $(x, y)$ とその最小値を求めよ。
  \item $0\leqq\theta<2\pi$ のとき，関数 $f(\theta)=2\sin 2\theta+2\sqrt{2}(\sin\theta+\cos\theta)$ について，$t=\sin\theta+\cos\theta$ とおくことで $f(\theta)$ を $t$ の式で表し，$t$ のとり得る値の範囲を求めよ。さらに，$f(\theta)$ の最大値と最小値，およびそのときの $\theta$ の値をそれぞれ求めよ。
  \item 次の極限値をそれぞれ求めよ。
  \[
    \text{(i)}\ \ \lim_{n\to\infty}\left(\sqrt{9n^{2}+6n+2}-3n\right)
    \qquad
    \text{(ii)}\ \ \lim_{x\to 0}\frac{(e^{2x}-1)\sin 3x}{1-\cos 2x}
  \]
  \item 等式 $\dsp f(x)=2\sin x+\frac{2}{\pi}\int_{0}^{\pi}f(t)\sin(x-t)\,dt$ を満たす連続関数 $f(x)$ を求めよ。
\end{shomon}
\end{daimon}

\begin{daimon}{2}
数列 $\{a_{n}\}$ が
\[
  a_{1} = 1,\qquad
  a_{n+1} = \frac{n}{n+2}a_{n} + \frac{3n+2}{(n+1)(n+2)} \quad (n=1, 2, 3, \ldots)
\]
によって定められている。次の問いに答えよ。\hfill（配点 20点）
\begin{shomon}
  \item $a_{2}$，$a_{3}$ の値をそれぞれ求めよ。また，$b_{n}=n(n+1)a_{n}$ とおくとき，$b_{n+1}-b_{n}$ を $n$ の式で表せ。
  \item 数列 $\{b_{n}\}$ の一般項，および数列 $\{a_{n}\}$ の一般項をそれぞれ $n$ の式で表せ。
  \item 極限値 $L=\lim\limits_{n\to\infty}a_{n}$ を求めよ。また，無限級数 $\dsp\sum_{n=1}^{\infty}\frac{L-a_{n}}{n}$ の収束・発散を調べ，収束するときはその和を求めよ。
  \item 無限級数 $\dsp\sum_{n=1}^{\infty}\frac{a_{n}}{n(n+2)}$ の和を求めよ。
\end{shomon}
\end{daimon}

\begin{daimon}{3}
複素数平面上で，$z\neq 2$ を満たす複素数 $z$ に対して
\[
  w = \frac{2z-1}{z-2}
\]
とおく。ただし，$i$ は虚数単位とする。次の問いに答えよ。\hfill（配点 20点）
\begin{shomon}
  \item $w\neq 2$ であることを示し，$z$ を $w$ の式で表せ。また，点 $z$ が単位円 $|z|=1$ 上を動くとき，点 $w$ が描く図形を求めよ。
  \item 点 $z$ が虚軸全体を動くとき，点 $w$ が描く図形を求め，複素数平面上に図示せよ。
  \item 点 $z$ が虚軸全体を動くとき，$|w|$ の最小値と，そのときの $z$ の値を求めよ。また，$|w|$ がとり得る値の範囲を求めよ。
  \item 点 $z$ が虚軸の負の部分（$z=ti$，$t<0$）を動くとき，$w$ の偏角を $\theta$（$-\pi<\theta\leqq\pi$）とする。$\tan\theta$ の最大値と，そのときの $w$ および $z$ の値を求めよ。
\end{shomon}
\end{daimon}

\clearpage

\begin{daimon}{4}
座標空間内に，原点 $\Pt{O}(0, 0, 0)$ を中心とする半径 $3$ の球面 $S: x^{2}+y^{2}+z^{2}=9$ と，$3$ 点 $\Pt{A}(3, 0, 0)$，$\Pt{B}(0, 3, 0)$，$\Pt{C}(0, 0, 3)$ を通る平面 $\alpha$ がある。次の問いに答えよ。\hfill（配点 20点）
\begin{shomon}
  \item 平面 $\alpha$ の方程式を求めよ。また，原点 $\Pt{O}$ から平面 $\alpha$ に下ろした垂線の足 $\Pt{H}$ の座標と，線分 $\Pt{OH}$ の長さを求めよ。
  \item 球面 $S$ と平面 $\alpha$ が交わってできる円を $K$ とする。円 $K$ の半径 $r$ を求めよ。また，点 $\Pt{P}$ が円 $K$ の周上を動くとき，内積 $\Vec{OA}\cdot\Vec{OP}$ の最大値と最小値をそれぞれ求めよ。
  \item 円 $K$ の周上の点 $\Pt{D}(-1, 2, 2)$ について，$\Pt{D}$ が円 $K$ 上にあることを確かめ，三角形 $\Pt{BCD}$ の面積を求めよ。
  \item 点 $\Pt{Q}$ が球面 $S$ 上を動くとき，四面体 $\Pt{BCDQ}$ の体積の最大値と，そのときの点 $\Pt{Q}$ の座標を求めよ。
\end{shomon}
\end{daimon}

\begin{daimon}{5}
$x>0$ で定義された関数 $f(x)=\dfrac{\log x}{x}$ のグラフを $C$ とする。ただし，対数は自然対数とし，$e$ は自然対数の底とする。次の問いに答えよ。\hfill（配点 20点）
\begin{shomon}
  \item 関数 $f(x)$ の増減・極値・凹凸・変曲点を調べ，さらに $\lim\limits_{x\to +0}f(x)$ および $\lim\limits_{x\to\infty}f(x)$ を求めて，曲線 $C$ の概形をかけ。
  \item 原点 $\Pt{O}(0, 0)$ から曲線 $C$ に引いた接線を $\ell$ とする。接線 $\ell$ の方程式と接点 $\Pt{T}$ の座標を求めよ。また，$x>0$ において不等式 $\dfrac{\log x}{x}\leqq\dfrac{x}{2e}$ が成り立つことを示せ。
  \item 曲線 $C$ と接線 $\ell$，および $x$ 軸で囲まれた部分の面積 $S$ を求めよ。
  \item 曲線 $C$ の $1\leqq x\leqq e$ の部分と $x$ 軸，および直線 $x=e$ で囲まれた部分を $x$ 軸のまわりに $1$ 回転してできる立体の体積 $V$ を求めよ。
\end{shomon}
\end{daimon}

\scratchpage
\end{monsheet}

